% SPDX-FileCopyrightText: 2026 Hasan Melih Akbulut
% SPDX-License-Identifier: CC-BY-4.0
\documentclass[11pt,a4paper]{article}

\usepackage[T1]{fontenc}
\usepackage[utf8]{inputenc}
\usepackage{lmodern}
\usepackage[margin=1in]{geometry}
\usepackage{booktabs}
\usepackage{array}
\usepackage{amsmath}
\usepackage{siunitx}
\usepackage{microtype}
\usepackage{tabularx}
\usepackage[hidelinks]{hyperref}   % loaded last, as it patches others

\sisetup{group-separator={,}, group-minimum-digits=4,
         group-digits=integer, text-series-to-math=true}

\newcolumntype{L}[1]{>{\raggedright\arraybackslash}p{#1}}
\newcommand{\file}[1]{\texttt{\small #1}}

\title{What the Deck Checked:\\
The Design-Side Residual of an Open-PDK Sign-off Run}

\author{Hasan Melih Akbulut\thanks{Independent researcher.
Correspondence: \texttt{57303760+Melihakbulut221@users.noreply.github.com}.
Sources and records are open; see \S\ref{sec:artefacts}.}}

\date{14 September 2026}

\begin{document}
\maketitle

\begin{abstract}
``DRC clean'' names one outcome and hides four populations: the rules a
shipped deck contains, the subset a run evaluates, the subset a design
actually fires, and where the resulting markers land. We account for all
four on one completed design---a routed \SI{130}{\nano\meter}
system-on-chip on the IHP SG13G2 open PDK containing eight vendor SRAM
macros---and report the residual.

The run's deck declares \num{173} rule categories in 47 groups. \textbf{Three
fire.} Every marker they produce lies inside the vendor macro hierarchy:
across three layouts of two different sizes, the count of markers naming
a cell that is not vendor-supplied is \textbf{zero}, on \num{192} and then
\num{204} distinct cells. The same PDK commit ships a second deck of
\num{371} further rules which the sign-off flow does not run, and which a
prior measurement found returns zero on this vendor macro.

Marker counts are not violation counts and the gap is stable rather than
incidental: two rules over one derived layer report an identical
(cell, geometry) set, so \num{11048} markers are \num{6584} distinct
shapes, a ratio of \num{1.68} that reproduces at \num{9668}/\num{5750}
on a six-macro layout of the same design.

The sharpest result is a distance. Running Magic's deck against the
LEF abstraction rather than the vendor GDS produces \num{182} error boxes
of which \num{20} fall inside a macro footprint; \textbf{every one of the
other \num{162} lies within \SI{0.500}{\micro\meter} of a macro edge}, in
six discrete values and no others, with nothing in the open field of
standard cells and routing. The bound reproduces at
\SI{0.500}{\micro\meter} on a second layout and, run backwards, at
\SI{0.380}{\micro\meter} on a six-macro layout measured eleven days
earlier---so the property held for a published measurement all along and
was not asked of it.

We also report what the deck is the \emph{only} instrument for: the
detailed router declines \texttt{LEF58\_ENCLOSURE} on all seven cut
layers this PDK defines, because the technology LEF states plain
\texttt{ENCLOSURE} and qualifies none of it with \texttt{CUTCLASS}. And
we report an asymmetry in the accompanying LVS: the comparison passes
only when the vendor macro is black-boxed on both sides, and the passing
comparison is \emph{weaker} than the one that fails, whose failure is a
bus-delimiter naming artefact rather than a connectivity error.

No silicon exists. Nothing here was fabricated or exposed to a beam.
Every number is a deck output on a layout.
\end{abstract}

\section{Introduction}
\label{sec:intro}

A design that has passed design-rule checking is usually reported as a
single fact. The report is a number and the number is zero, or the
number is not zero and the design is not finished. Neither form says
which rules were in the deck, which of them the run evaluated, which of
those the design contained geometry capable of firing, or---when a
marker does appear---whether it names something the design drew or
something it instantiated.

Those four questions have different answers and only the last is
commonly asked. This paper asks all four of one design and reports the
residual: what a shipped open PDK's sign-off deck actually checked, and
what it never exercised.

The object is a routed \SI{130}{\nano\meter} system-on-chip on IHP
SG13G2, an open process design kit~\cite{ihppdk}, hardened through an entirely open
flow. It contains \num{168737} instances and eight vendor SRAM macros of
three types. It is not a small design and it is not a clean one: it
misses its timing constraint at the slow corner by
\SI{7.35}{\nano\second}. That is stated here because a deck result is
usually presented alongside a design that closed, and the accounting
below does not depend on closure.

\paragraph{What this is not.} It is not a claim that the deck is wrong.
It is not a radiation result, a silicon result, or a comparison of open
tools against commercial ones. It is an accounting of one run.

\paragraph{Contributions.}
\begin{enumerate}
\item A four-population accounting of one sign-off run: deck contents,
  run coverage, design firings, and marker attribution
  (\S\ref{sec:accounting}).
\item The attribution result, at two design sizes and on three layouts:
  zero markers name a cell the design drew (\S\ref{sec:attribution}).
\item A distance bound on the abstraction's artefacts that replaces a
  mechanistic argument with a measurement, including the backwards test
  that shows it held for an earlier published run (\S\ref{sec:distance}).
\item The complementary negative: the rules the router declines, which
  makes the deck the only instrument looking at them
  (\S\ref{sec:router}).
\item An asymmetry in the accompanying LVS, in which the passing
  configuration checks strictly less than the failing one
  (\S\ref{sec:lvs}).
\end{enumerate}

\section{The four populations}
\label{sec:accounting}

\subsection{What the deck contains}

The PDK ships its KLayout rules as a tree of \num{39} files under
\file{libs.tech/klayout/tech/drc/rule\_decks/}, of which \num{38}
contain at least one rule. Counting \texttt{.output()} calls statically
across them gives \num{376} distinct rule names, of which \num{25} are
constructed by string interpolation over a layer index---so
\texttt{M\#\{met\_no\}.a} is one call and several rules.

That count divides unevenly. \textbf{Three hundred and fifty-one of the
\num{376} are in a single file}, \file{sg13g2\_maximal.drc}, which the
sign-off flow does not run. The PDK's own README distinguishes the two:
the main set is ``intended for foundry precheck purposes'' and its rules
``have been thoroughly verified, tested, and optimized''; the maximal set
is ``additional residual rules \ldots{} not been verified or tested and
may be slower''. The PDK's own \file{run\_drc.py} runs both by default.
The flow used here runs only the first.

A prior measurement on this PDK~\cite{companion} ran the maximal deck
deliberately: it
evaluates \num{371} rules and returns zero, both on the vendor macro
standalone and on a smaller full design, and a positive control built
from four shapes with two deliberate violations shows it reporting five
errors---so the zeros are results and not a silent no-op. That
measurement is not repeated here and is cited rather than re-derived.

\subsection{What the run evaluated}

The report the flow emits is a KLayout \file{.lyrdb}, and it declares a
category for every rule the run registered. On this design that is
\textbf{\num{173} categories across \num{47} groups}: \texttt{M1} through
\texttt{M5}, \texttt{V1} through \texttt{V4}, \texttt{TM1}, \texttt{TM2},
\texttt{MIM}, \texttt{NW}, \texttt{Act}, \texttt{Gat}, \texttt{Seal},
\texttt{Pad}, the \texttt{Fil} family, and the rest.

The \num{173} is the number that matters for reading the result, and it
is not printed anywhere in the flow's own summary. A reader is given a
count of errors. Whether that count came from a deck of three rules or
three hundred is not in the artefact the flow presents.

\subsection{What the design fired}

\textbf{Three categories fire.} They are the same three on every layout
of this design we measured:

\begin{center}
\begin{tabular}{lrrr}
\toprule
& one macro & six macros & eight macros \\
\midrule
\texttt{Sdiod.d}       & --- & \num{3918} & \num{4464} \\
\texttt{Sdiod.e}       & --- & \num{3918} & \num{4464} \\
\texttt{Cnt.c.digibnd} & --- & \num{1832} & \num{2120} \\
\midrule
markers                & \num{2316} & \num{9668} & \num{11048} \\
distinct shapes        & \num{1294} & \num{5750} & \num{6584} \\
ratio                  & \num{1.79} & \num{1.68} & \num{1.68} \\
distinct cells         & \num{57}   & \num{192}  & \num{204} \\
\bottomrule
\end{tabular}
\end{center}

The one-macro column is a prior measurement on the vendor macro alone and
is included for the series; its per-rule split is in that record.

\paragraph{Markers are not violations, and the gap is structural.}
\texttt{Sdiod.d} and \texttt{Sdiod.e} are evaluated over the same derived
layer. Comparing their items as (cell, geometry) pairs, the intersection
is the whole of both: \num{4464} shared, \num{0} in either alone. Every
flagged shape is reported twice. So \num{11048} markers are \num{6584}
distinct shapes, and the ratio \num{1.68} reproduces exactly on the
six-macro layout. A headline that does not say which it is has a
\num{68}\,\% error bar attached to it and no indication that it does.

\section{Where the markers are}
\label{sec:attribution}

Each item in the report names the cell it lies in. That makes
attribution a lookup rather than a geometric test, which matters: a
containment test against a footprint cannot distinguish a marker
\emph{in} the macro from a marker the macro's presence provoked
\emph{beside} it, and the hierarchy name can.

\begin{center}
\begin{tabular}{lrrr}
\toprule
layout & distinct cells & vendor cells & \textbf{cells the design drew} \\
\midrule
six macros, \file{s77gate}   & \num{192} & \num{192} & \textbf{0} \\
eight macros                 & \num{204} & \num{204} & \textbf{0} \\
eight, antenna-repaired      & \num{204} & \num{204} & \textbf{0} \\
\bottomrule
\end{tabular}
\end{center}

Every cell named by every marker is an \texttt{RM\_IHPSG13\_1P\_ROWDEC*},
\texttt{COLDEC*} or \texttt{RSC\_IHPSG13\_*}---row decoders, column
decoders and the standard cells inside the vendor's own array. The
design's standard cells, its routing, its power grid and its pads
produce nothing.

\paragraph{The third row is a measurement and not a repetition.} The
antenna-repaired layout differs from the second by \num{60} inserted
diodes and a re-route. Every figure above is identical: the same
\num{11048}, the same \num{6584}, the same \num{204} cells, the same
equality between the two Schottky rules. Sixty cells of the design's own
geometry entered the layout and the deck's output did not move. That is
the control this section would otherwise lack.

\paragraph{Scale, and what a sweep settles that a series cannot.} From
six macros to eight, markers rise by \num{1380} and cells by \num{12};
the added markers name cell types that exist only inside the added
macro, and no pre-existing cell's count changes.

The three-point series \num{2316}, \num{9668}, \num{11048} moves array
size and instance count together and cannot say which is responsible. So
all \num{27} macros the PDK ships were run standalone through the same
deck. \textbf{Width does not move the count at all}: \texttt{1024x8},
\texttt{1024x16} and \texttt{1024x32} return an identical \num{2672}
across a fourfold difference in stored bits, and \texttt{256x16},
\texttt{256x32} and \texttt{256x48} an identical \num{2148}. Width
enters in exactly one place and by the same amount every time:
\texttt{x64} adds \num{206} at each of three depths---\num{2354}
against \num{2148}, \num{2522} against \num{2316}, \num{2878} against
\num{2672}. \textbf{Depth does move it}: \num{2148}, \num{2316},
\num{2672}, \num{3538} for \num{256} through \num{2048}. A second
port roughly doubles it.

The population is therefore generated by the decoder and port structure
the vendor draws, not by the memory the design asked for---a design that
halves its word width buys no reduction whatever. \texttt{512x16}
returns \num{2316}, anchoring the sweep to the one-macro measurement.

\section{A distance, where an argument used to be}
\label{sec:distance}

Magic's deck can be run two ways: over the vendor GDS, or over the LEF
abstraction, in which each macro is the obstruction rectangles its
abstract declares. The first mode on the six-macro layout produces
\num{39969214} error boxes in \SI{17.2}{\hour}, all of them inside a macro
footprint. The second is cheap and is the one reported here.

Abstracted, the eight-macro layout produces \num{182} boxes: \num{20}
inside a macro footprint, \num{162} outside. The standard reading of the
\num{162} is mechanistic---a LEF \texttt{OBS} blanket is one shape wider
than \SI{10}{\micro\meter}, so ordinary routing beside it violates a
wide-metal spacing rule it would not violate against the macro's real
geometry. That reading is an interpretation. It can be replaced.

For each box outside a footprint, take the distance to the nearest of the
eight rectangles the placement reserves:

\begin{center}
\begin{tabular}{lrrr}
\toprule
gap to nearest macro edge & \multicolumn{3}{c}{boxes} \\
\cmidrule(l){2-4}
 & eight macros & eight, repaired & six macros \\
\midrule
\SI{0}{\micro\meter} (touching) & 2  & 2  & 4  \\
\SI{0.220}{\micro\meter}        & 31 & 31 & 1  \\
\SI{0.260}{\micro\meter}        & 39 & 55 & 27 \\
\SI{0.280}{\micro\meter}        & 4  & 5  & 21 \\
\SI{0.320}{\micro\meter}        & 54 & 54 & 51 \\
\SI{0.380}{\micro\meter}        & 15 & 12 & 3  \\
\SI{0.500}{\micro\meter}        & 5  & 3  & 0  \\
\midrule
total outside                   & \num{150} & \num{162} & \num{107} \\
\textbf{maximum}                & \textbf{\SI{0.500}{\micro\meter}}
                                & \textbf{\SI{0.500}{\micro\meter}}
                                & \textbf{\SI{0.380}{\micro\meter}} \\
\bottomrule
\end{tabular}
\end{center}

The gaps take six values and no others. Not one box sits in the field of
standard cells and routing where almost every polygon in this layout is.
The claim the mechanism was advanced to support---that these are
artefacts of the boundary the abstraction invents---is now a distance
rather than an interpretation, and a reader can check it with the
placement and the report.

\paragraph{It held for a published run and was not asked of it.} The
third column is a six-macro layout measured eleven days before this
analysis existed, whose \num{107} outside boxes had been reported without
a distance. Its maximum is \SI{0.380}{\micro\meter}. The property was
true of that measurement all along. This is the ordinary way a result is
missed: not by measuring wrongly, but by not putting a question to an
instrument that could already answer it.

\paragraph{One caution, stated because it changes a number.} A second
split is available---against the drawn extent, the LEF box plus a
\SI{0.225}{\micro\meter} well overhang---and it should be read with care.
Of the \num{33} boxes it reclassifies as straddling an edge, \num{31}
overlap by exactly \SI{0.005}{\micro\meter}, at three different macro
edges. The markers sit \SI{0.220}{\micro\meter} from the LEF edge and the
\SI{0.225}{\micro\meter} model overshoots them by five nanometres. Set
the constant to \SI{0.220}{\micro\meter} and all \num{31} move. That
split is decided by a modelling choice below the resolution of the thing
it models; the footprint split above it is not.

\section{What the deck is the only instrument for}
\label{sec:router}

An accounting of what a deck checks is incomplete without the
complementary question: what checks nothing else does.

Every detailed-routing log this design has produced carries seven
instances of
\texttt{[DRT-0349] LEF58\_ENCLOSURE with no CUTCLASS is not supported.
Skipping for layer X}---for \texttt{Cont}, \texttt{Via1} through
\texttt{Via4}, \texttt{TopVia1} and \texttt{TopVia2}, which is every cut
layer the technology defines.

The cause is in the technology LEF, and it is not a defect: the file
states plain LEF \texttt{ENCLOSURE} on all seven cut layers and qualifies
none of them with \texttt{CUTCLASS}. The router reads those as
\texttt{LEF58\_ENCLOSURE} rules it cannot interpret and declines all
seven.

So for via enclosure on this PDK, the sign-off deck is not a second
opinion. It is the only opinion. A design that runs with the deck
switched off---which is a supported configuration, and was this project's
configuration for an extended period---has no instrument checking cut
enclosure at all. That is a stronger statement about the value of the
\num{173} than the firing count alone supports, and it runs in the
opposite direction.

\section{The LVS asymmetry}
\label{sec:lvs}

The same run tree carries a layout-versus-schematic comparison, and it
has a property worth reporting alongside the deck: \emph{the
configuration that passes checks strictly less than the configuration
that fails}.

Supplying the vendor CDL, netgen expands the macro's transistors on the
schematic side while the extraction sees the LEF abstract on the layout
side. The comparison fails: \num{64901} netlist nets against \num{63155}
layout nets, and a top-level pin-matching failure. The cause is
\emph{naming}. The CDL spells buses \texttt{A\_DIN<32>} and globals
\texttt{VSS!}, \texttt{VDD!} where the layout has \texttt{A\_DIN[32]},
\texttt{VSS}, \texttt{VDD}; the net difference follows from the
delimiters.

Black-boxing the macro on both sides, the comparison passes:
\num{63155} = \num{63155}, circuits match uniquely, all seven difference
counters zero. But circuit~2 is then built with no SRAM model at all, the
extracted side carries an empty subcircuit per macro, and \num{974} macro
pins are disconnected nodes on both sides alike. What passes is top-level
connectivity into the macro pins. Nothing inside the macro is compared,
and nothing was going to be: the vendor CDL covers two of the three macro
types instantiated here, so the third was a black box in the failing run
as well.

Dropping the CDL does not repair the comparison. It removes the part of
it that looked inside the macro. Both facts belong in a report of what
the run checked.

\section{An error class that only the larger design has}
\label{sec:m2d}

Of the \num{20} Magic boxes inside a macro footprint, four are a rule the
six-macro layout cannot produce: \texttt{M2.d}, minimum area, twice in
each of the two \SI{512}{}$\times$\SI{16}{} check macros, at
\SI{0.26}{\micro\meter} square.

They are the macro's own pin rectangles for \texttt{A\_DOUT[14]} and
\texttt{A\_DOUT[15]}. One hundred and eight of the macro's \num{111} pins
are exactly that size; those two are the ones with nothing routed to
them, so no attached wire grows the shape past the minimum. The design
uses fourteen of the sixteen bits---two words' worth of seven check bits
each---and the top two outputs are unconnected, which the routed netlist
confirms: those two nets appear once, at the macro port, where every used
bit appears twice.

In the macro's real geometry the pin attaches to internal metal and has
ample area. The marker exists because the run reads the abstract. It is
reported here because it is an instance of the general point with a
concrete mechanism: the error population a deck produces is a function of
the abstraction the run was given, and it changes when the design grows
in a way that has nothing to do with the design being more or less
correct.

\section{Threats to validity}
\label{sec:threats}

\paragraph{One design, one PDK.} Everything here is one system-on-chip on
one process. The attribution result reproduces across three layouts and
two design sizes, but those are three layouts of the same design.

\paragraph{The expensive mode was not repeated.} The
\SI{17.2}{\hour} full-hierarchy Magic run exists for the six-macro layout
only. Neither the \num{39969214} nor the finding derived from it has an
eight-macro counterpart, and the two macro types added since have never
been checked standalone.

\paragraph{The deck is not validated here.} We report what it did on this
design. Whether its rules are correct is a different question with its
own literature, and the PDK's own issue tracker carries open items on it.

\paragraph{The design does not close timing.} It misses the slow corner
by \SI{7.35}{\nano\second}. Nothing in the accounting depends on that,
but a reader should not take this design as an example of a finished one.

\section{Related work}
\label{sec:related}

Krinke, Fischbach and Lienig~\cite{krinke2024} measured open-source
layout verification against a commercial PDK and characterised what a generated deck can
express. That work is tool-side: it asks what the checker is capable of.
This paper is design-side: it asks what one run, of a shipped deck, on
one completed layout, actually exercised. The two compose and neither
substitutes.

The vendor-macro exclusion is documented in open-silicon practice as a
workaround---precheck flows carry flags to exclude SRAM from
checks---and the reading here is not that the exclusion is novel but that
its \emph{extent} on a real design has not been reported: how many rules
remain, which fire, and whether anything the design drew is among them.

The prior measurement of the second deck on this PDK, its positive
control, and the marker-versus-shape correction that this paper's
\num{1.68} ratio extends, are the author's own~\cite{companion} and are
cited as such rather than rederived.

\paragraph{The author's own nearest neighbour.} ZIRH-2~\cite{zirh2} is a
rad-hard
experiment chip on this same open PDK, public under Apache-2.0, carrying
triplicated logic and error-coded memory through an open flow to a routed
layout. It is named here because it is the closest artefact to this one
and because a paper that does not cite its author's adjacent work invites
the reading that it was not looked for. It does not carry the accounting
this paper reports.

\section{Artefacts}
\label{sec:artefacts}

The classifier that produces the footprint split and the distances is
\file{scripts/classify\_magic\_drc.py}. It reproduces a previously
recorded classification of the six-macro layout exactly---\num{16} inside,
\num{0} straddling, \num{107} outside of \num{123}---which is the check
that it is measuring what the earlier, uncommitted script measured.

Two defects in it are recorded because they shaped the numbers. It
refuses to report when it finds no macro in the placement, after a first
version anchored its pattern at the wrong column, found none, and
reported every box as ``outside''---the most useful-looking answer
available and the exact inverse of the result. And a macro edge lands on
a binary-float sum, so a box touching an edge over zero area counted as
straddling it until a tolerance was added; that was one box in \num{123},
and it was the difference between reproducing the recorded number and
quietly disagreeing with it.

\bibliographystyle{plain}
\bibliography{refs}

\end{document}
